% ---------------------------------------------------------------------------
% Section 1 -- Introduction (short paper: FLAT, no subsections)
%
% Author's text, compressed from 1,230 to ~1,000 words. Nothing was cut from
% the argument -- the reductions are merges and removed restatement:
%   * P3 and P4 merged (the "In all these cases" opener restated P3's close)
%   * P7 and P8 merged (vulnerability scoring folded into the same point)
%   * the three-factor paragraph keeps its three claims, loses two of three
%     illustrative lists (the payment/path/scope list and the transfer example,
%     which paragraph 3 already carries)
%   * "We therefore introduce" and "We implement this" merged
%   * contributions tightened; all three claims intact
%
% Author's softenings preserved: "Much recent research", "has accompanied the
% growth", "not by itself sufficient", "do not, however, express".
% ---------------------------------------------------------------------------
\section{Introduction}
\label{sec:intro}

AI agents no longer only generate text; they act. They read data, modify files,
update databases, and call external services for their users. The Model Context
Protocol (MCP)~\cite{anthropic2024mcp} standardizes this exchange: a server
publishes tools, and a connected agent discovers and invokes them over
JSON-RPC~\cite{mcp2026spec}. The public server ecosystem is growing
quickly~\cite{hou2026mcp}, and a study of 1{,}899 open-source servers found
recurring security weaknesses~\cite{hasan2025firstglance}. An agent's decision
now lands directly on the systems and resources behind those tools. MCP does
carry a risk vocabulary, but a coarse one: a tool descriptor may declare
\texttt{readOnlyHint}, \texttt{destructiveHint}, \texttt{idempotentHint} and
\texttt{openWorldHint}~\cite{mcp2026spec}. These are booleans, declared per tool
rather than per call, and asserted by the server about itself; the specification
tells clients to treat them as untrusted unless the server is trusted. Its
remaining answer to risk is a human who can deny the invocation. A hint labels a
capability. It says nothing about which resource a given call
reaches~\cite{hou2026mcp,buhler2026agentbound}.

Most MCP and agent security work studies the agent's side: how a malicious or
compromised server subverts the agent connecting to
it~\cite{invariant2025toolpoisoning,shi2025toolhijacker}. A smaller line of work
defends the server. Before deployment, static analysis maps implementation
weaknesses to attack patterns and scores a repository~\cite{kumar2026mcpinsos},
or audits declared capabilities for over-privilege~\cite{huang2026auditing}. At
runtime, access control has been expressed as a capability manifest enforced by
a sandbox~\cite{buhler2026agentbound}, and hybrid analysis pairs a call's
parameters with its observed execution~\cite{he2026mtguard}. Where these methods
score, the unit is a server or a capability, never an invocation. They also
defend the opposite direction: a capability sandbox confines a faulty or
malicious server to protect the host~\cite{buhler2026agentbound}, while the
assets at stake here are the ones the server itself owns.

We study \emph{misuse}: the inappropriate use of a legitimate capability. Misuse
needs no malicious agent and no malicious server. It starts with an operator who
states a task loosely, delegates too broadly, or grants more permission than the
task needs; measurements of MCP servers report broad system access and
over-privileged tools~\cite{li2025privilege,huang2026auditing}. It starts with
the model, which can pick an available but wrong tool~\cite{yin2025reasoningtrap},
or know a risk in the abstract and still miss it in a live
trace~\cite{tang2025riskknowledge}. It needs no single harmful call: legitimate
tools chain into a workflow that is harmful in combination, driven by
instructions hidden in data the agent reads during an ordinary
task~\cite{zhao2026parasites}. Prompt injection~\cite{zhan2024injecagent} and
agentic misalignment~\cite{lynch2025misalignment} reach the same place by other
routes, but are not our subject. We do not ask why the agent issued the call. We
ask what the call would do. The same \texttt{read\_file} is routine against a
build log and a disclosure against a credentials store. Our unit of analysis is
the whole invocation: the tool, the resource it reaches, the arguments, the
execution context, and the agent's recent activity.

Access control decides whether a principal may reach a resource. Role- and
attribute-based models~\cite{sandhu1996rbac,hu2014abac} decide whether an
operation is permitted; they attach no magnitude to a permitted one. MCP widens
the gap: servers hold authorization as persistent state instead of re-checking
each call~\cite{huang2026caller}, and agent frameworks expose a capability
without re-authorizing the concrete call it emits~\cite{mellafe2026capability}.
Authentication says who is calling, authorization says what they may do. Neither
says what this call costs.

Adding a computed risk value to that decision is not new. Quantified
risk-adaptive access control appeared in 2007~\cite{cheng2007fuzzymls} and was
folded into attribute-based models in 2011~\cite{kandala2011radac}. The family
is surveyed~\cite{atlam2020riskbased,cheimonidis2023dynamic}, deployed where a
numeric trust value gates access~\cite{bradatsch2024zerotrust}, and recently
extended to agents, where an {LLM} judge weighs target-resource risk against its
own uncertainty~\cite{fleming2025tbac}. These models score a subject--object
pair: this principal, this resource, this context. Two calls to the same tool by
the same agent score alike.

No existing method estimates what a concrete MCP invocation would cost. Most MCP
and agent defenses return an enforcement decision, allow or
block~\cite{xing2025mcpguard,jing2025mcip,shi2025progent,zhu2025miniscope,abaev2026agentguardian}.
Recent systems return richer outcomes --- consent levels, verdict classes,
calibrated detector
scores~\cite{li2026conleash,yang2026agenttrust,zhang2026stars,hossain2026nexus}
--- but a containment level, an action class and a detection confidence are not
magnitudes. Others validate calls at runtime~\cite{betser2026agentrim}, learn
access-control policies~\cite{abaev2026agentguardian}, infer risk from
logs~\cite{fu2025riskcue}, or check arguments against a declared
intent~\cite{zhu2026igac,rahman2026rolestratified}; none prices the resource the
call reaches. Vulnerability scoring does not either: CVSS and AIVSS rate a
vulnerability or a system, not a live
call~\cite{spring2021time,bahar2024validity,owasp2025aivss}. A magnitude matters
because it separates calls to allow, calls to review, and calls to
refuse~\cite{cao2024mad}.

Risk belongs to a (tool, resource) pair, not to the tool: \texttt{read\_file} is
low risk on a public document and high risk on regulated data. A baseline must
combine tool impact, resource sensitivity, and \emph{blast radius} --- how far
one call reaches. Arguments move that baseline, and the same call reads
differently against different histories.

We compute a graded risk score at two times. At design time we score every
(tool, resource) pair a server exposes, from tool impact, resource sensitivity,
and blast radius. Sensitivity comes from the organization's own
information-classification policy, not from a number assigned without
one~\cite{nist2004fips199,nist2008sp80060}. At request time we adjust that
baseline with the actual arguments, the execution context, and the agent's
recent activity. The score is compared against organizational thresholds, so a
call is approved, sent for human review, or refused (\S\ref{sec:framework}).
Deciding when an agent proceeds alone and when control passes to a person is the
central question of adjustable autonomy~\cite{pynadath2002adjustable}; the score
makes that threshold computable. We implement this as an enforcement framework
and evaluate it on tool catalogs captured from real vendor MCP servers, paired
with organizational policies, resource inventories, and invocation scenarios
(\S\ref{sec:eval}).

This paper makes three contributions. First, a \textbf{risk-analysis methodology}
for assessing the misuse risk of an MCP invocation from tool impact, resource
sensitivity, blast radius, runtime arguments, execution context, and the agent's
recent activity, deriving resource sensitivity from organizational policy rather
than from a fully labeled asset inventory. Second, an \textbf{enforcement
framework} that combines the design-time and request-time assessments and uses
them to make risk-based access-control decisions for MCP servers. Third, an
\textbf{evaluation setting} that pairs tool catalogs from real vendor MCP servers
with synthesized organizational policies, resource inventories, and invocation
scenarios, used to measure whether the calculated risk levels match
organizational judgment and whether runtime arguments and behavioral history
affect the score as intended.

% ---------------------------------------------------------------------------
% IF IT MUST GO BELOW 900 -- the next cuts all cost something:
%   * The three-factor paragraph (~95w) duplicates contribution 1. Cutting it
%     loses the read_file example, which is the only concrete illustration of
%     why risk is a (tool, resource) property.
%   * Paragraph 2 (~55w) can fold into paragraph 3, but the reverse-direction
%     statement is the paper's positioning and reads better standing alone.
%   * The graded-estimate sentence (~30w) can go, but it is the justification
%     for producing a score rather than a verdict.
%
% CONTRIBUTIONS AS A LIST -- six of six sampled AAMAS papers use one:
%   \begin{enumerate}[leftmargin=*,itemsep=1pt,topsep=2pt]
%     \item A \textbf{risk-analysis methodology} ...
%     \item An \textbf{enforcement framework} ...
%     \item An \textbf{evaluation setting} ...
%   \end{enumerate}
% ---------------------------------------------------------------------------
